\chapter{Fundamentação Teórica}
\label{cap:fundamentacao-teorica}

A estimação da potência ativa em corrente alternada (PCA) de inversores
fotovoltaicos a partir de variáveis meteorológicas combina conceitos de
conversão fotovoltaica, aprendizado supervisionado e inferência em
microcontroladores. Este capítulo organiza esses fundamentos na ordem em
que aparecem no restante do trabalho: primeiro, os fatores físicos que
condicionam a geração e a definição de uma estimativa contemporânea;
depois, o conceito de sensor virtual, a organização temporal dos dados e
as transformações aplicadas ao alvo; em seguida, os modelos de regressão
e as métricas de avaliação; por fim, os requisitos de execução embarcada.
Cada seção fixa as hipóteses e os compromissos que sustentam as escolhas
metodológicas descritas adiante.

\section{Fundamentos da geração fotovoltaica}
\label{sec:ft-geracao}

\subsection{Irradiância, temperatura e potência elétrica}
\label{subsec:ft-fv}

A conversão fotovoltaica produz energia em corrente contínua, que o
inversor converte em corrente alternada. A potência entregue depende,
principalmente, da irradiância incidente no plano dos módulos, da
temperatura das células e das perdas do sistema. Modelos físicos
descrevem essa dependência a partir das características elétricas dos
módulos \cite{desoto2006}, e uma aproximação usual para a potência em
corrente contínua é dada por \cite{dobos2014}:

\begin{equation}
\label{eq:ft-potencia-fv}
P_{\mathrm{CC}}\approx P_{\mathrm{STC}}
\frac{G_{\mathrm{POA}}}{G_{\mathrm{STC}}}
\left[1+\gamma\left(T_{\mathrm{cel}}-T_{\mathrm{STC}}\right)\right],
\end{equation}

\noindent em que $P_{\mathrm{STC}}$, $G_{\mathrm{STC}}$ e
$T_{\mathrm{STC}}$ são valores de referência, $G_{\mathrm{POA}}$ é a
irradiância no plano dos módulos, $T_{\mathrm{cel}}$ é a temperatura da
célula e $\gamma$ é o coeficiente térmico de potência, geralmente
negativo. O sinal do coeficiente explica por que o aquecimento reduz a
potência mesmo quando a irradiância permanece alta. A irradiância
efetiva é medida no plano dos módulos; quando o dado disponível está no
plano horizontal, é necessário transpor a irradiância, e a escolha do
modelo de transposição afeta a precisão de modelos fotovoltaicos
\cite{dolara2015modelos,mayer2020modelos}. Neste trabalho, a variável de entrada já
corresponde à irradiância transposta, o que remove essa etapa do ajuste,
mas não elimina a incerteza associada à medição e à orientação dos
sensores.

A temperatura das células não é medida diretamente em muitas
instalações. Modelos empíricos estimam essa temperatura a partir da
temperatura ambiente, da irradiância e, em alguns casos, do vento e da
forma de montagem dos módulos. O modelo de Jacques et al.
\cite{jacques2013} segue essa abordagem e fornece uma estimativa de
temperatura do módulo em regime estacionário. O uso dessa estimativa
como entrada evita a instalação de sensores de temperatura na célula,
mas introduz a incerteza do próprio modelo e exige que as condições de
treinamento sejam representativas das condições de operação. A
temperatura ambiente, portanto, não substitui a temperatura do módulo:
ela é um dos argumentos do modelo que a estima.

A potência CA depende ainda do rendimento de conversão e do limite de
saída do inversor. Em inversores conectados à rede, a potência CA cresce
com a potência CC até o limite nominal; acima desse ponto, o
equipamento satura a saída e parte da energia disponível é perdida
\cite{good2016}. O fenômeno, conhecido como saturação ou
\textit{clipping}, introduz uma
não linearidade adicional na relação entre irradiância e PCA e explica
por que observações na vizinhança do limite nominal requerem tratamento
próprio na preparação dos dados. Além do limite de potência, o
rendimento varia com a carga e com a tensão de entrada, o que reforça a
diferença entre a potência CC estimada pela equação e a potência CA
medida.

Em conjunto, esses fatores mostram que a relação entre as condições
meteorológicas e a saída elétrica não é estritamente linear. Um modelo
puramente linear pode capturar a tendência média, mas tende a errar com
mais frequência nos extremos de irradiância e de temperatura, onde os
efeitos térmicos e a saturação do inversor se acumulam. Essa expectativa
orienta a comparação entre a regressão linear e os modelos não lineares
no capítulo de resultados.

\subsection{Estimação contemporânea da potência}
\label{subsec:ft-potencia-contemporanea}

Diferentemente da previsão, que estima uma grandeza em um instante
futuro, a estimação contemporânea determina a potência correspondente ao
mesmo intervalo das variáveis de entrada. A distinção importa porque a
previsão dispõe de informação defasada e precisa lidar com a evolução
futura das condições, enquanto a estimação contemporânea trabalha com
uma relação instantânea entre variáveis observadas no mesmo intervalo.
Quando a medição da PCA fica indisponível, esse tipo de estimador pode
produzir uma referência alternativa a partir de informações
meteorológicas que permaneçam acessíveis.

Em termos gerais, um estimador contemporâneo pode ser expresso como

\begin{equation}
\label{eq:ft-funcao-estimacao}
\widehat{P}_{\mathrm{CA}}(t)=f_{\boldsymbol{\theta}}(\mathbf{x}(t)),
\end{equation}

\noindent em que $\mathbf{x}(t)$ reúne as variáveis auxiliares
disponíveis no mesmo intervalo da potência estimada e
$\boldsymbol{\theta}$ representa os parâmetros aprendidos. A função
$f_{\boldsymbol{\theta}}$ pode ser linear ou não linear, e a escolha
depende da estrutura dos dados e do custo de implantação. A definição
operacional das entradas, a origem dos dados e a partição temporal são
especificadas na Metodologia. Como a estimativa depende da qualidade e
da continuidade dessas variáveis, a avaliação do estimador precisa
considerar não apenas o erro médio, mas também os regimes em que as
entradas permanecem disponíveis.

\section{Sensores virtuais e lacunas de te\-le\-me\-tri\-a}
\label{sec:ft-softsensor}

Um sensor virtual, ou \textit{soft-sensor}, estima uma variável de
interesse com base em medidas auxiliares e em um modelo computacional.
Em abordagens orientadas por dados, a relação entre entradas e saída é
aprendida durante o treinamento e reutilizada em novas observações
\cite{kadlec2009}. O sensor virtual não substitui o medidor físico; ele
fornece uma estimativa quando a medição direta está ausente ou é
insuficiente. Sua confiabilidade depende da precisão do modelo, da
qualidade das entradas e da estabilidade da relação aprendida ao longo
do tempo. No contexto fotovoltaico, sensores virtuais já foram
empregados no monitoramento em tempo real de sistemas de potência
apoiados por Internet das Coisas \cite{boubakr2022monitoramento}.

Interrupções de comunicação, aquisição ou armazenamento geram lacunas
nas séries de monitoramento. A literatura destaca a importância de
distinguir registros ausentes de medidas válidas de potência nula, bem
como de tratar inconsistências antes de avaliar o desempenho
fotovoltaico \cite{lindig2020monitoramento,livera2020qualidade}. Uma
leitura de potência igual a zero pode significar ausência de geração,
falha de medição ou um código de erro do equipamento; tratar todos os
zeros como observações válidas distorce a distribuição do alvo e a
avaliação do modelo. Por isso, a preparação dos dados combina filtros
físicos e a manutenção de zeros apenas nas faixas em que eles são
plausíveis.

A interpolação utiliza valores vizinhos da própria série e pode
requerer observações posteriores ao intervalo ausente. Em uma falha de
telemetria em tempo real, esse requisito nem sempre pode ser atendido,
porque a informação futura ainda não existe quando a estimativa é
necessária. Um \textit{soft-sensor} contemporâneo opera em outra lógica:
ele não reconstrói a série a partir dela mesma, mas calcula a potência a
partir de variáveis auxiliares medidas no mesmo instante. Essa diferença
é central para o problema tratado neste trabalho, em que a PCA é a
variável ausente e as variáveis meteorológicas continuam disponíveis.

A imputação de dados faltantes é estudada também no contexto da previsão
fotovoltaica \cite{lee2024}. Entretanto, lacunas artificiais inseridas em
registros históricos não reproduzem necessariamente todas as condições de
uma falha real, sobretudo quando se pressupõe a integridade dos sensores
meteorológicos. A avaliação com lacunas artificiais mede a resposta do
estimador em períodos escolhidos, mas não substitui a validação em campo,
com falhas simultâneas, atrasos de comunicação e deriva dos sensores.
Essa distinção limita o alcance dos resultados e é retomada no capítulo
de resultados.

\section{Representação dos dados e generalização}
\label{sec:ft-dados}

\subsection{Partição temporal e granularidade}
\label{subsec:ft-granularidade}

Na regressão supervisionada, cada observação associa um vetor de
entradas $\mathbf{x}_i$ à potência medida $y_i$. O treinamento ajusta os
parâmetros, a validação orienta a seleção de configurações e o teste
verifica o desempenho em observações não utilizadas no ajuste. Em séries
temporais, partições aleatórias podem aproximar excessivamente os
conjuntos e produzir métricas otimistas, porque observações vizinhas
compartilham tendências e condições meteorológicas. Por isso, a avaliação
deve respeitar a ordem temporal dos dados \cite{roberts2017}, com
treinamento no passado e teste no futuro.

A granularidade define o intervalo representado por cada registro. A
agregação por média suaviza oscilações rápidas da potência, mas também
modifica a distribuição do alvo e a quantidade de amostras
\cite{mayer2021}. Em resoluções
mais grossas, a variância do alvo diminui e a autocorrelação aumenta;
como consequência, as métricas de erro não são diretamente comparáveis
entre resoluções, e um MAE menor em 1~d não significa que o modelo seja
mais preciso em 5~min. Além disso, a média das medições diurnas não
equivale à energia diária nem à potência média durante 24 horas, porque
o período noturno e as lacunas de aquisição ficam fora do cálculo.

Essas propriedades justificam treinar um modelo por granularidade e
comparar os resultados dentro de cada resolução, em vez de eleger uma
resolução apenas pelo erro absoluto. A escolha da granularidade de
operação depende do compromisso entre a frequência do acompanhamento, o
volume de dados e o custo de treinamento e inferência.

\subsection{Assimetria, normalização e transformação de Box--Cox}
\label{subsec:ft-boxcox}

A distribuição da PCA pode concentrar observações em faixas de baixa
potência, reduzindo a representação relativa dos regimes menos
frequentes. Esse desequilíbrio influencia o ajuste de modelos baseados
em erros agregados: a perda é dominada pelas observações numerosas, e os
picos de potência recebem menos atenção durante o treinamento. Uma
distribuição mais simétrica, porém, não garante melhor predição; a
qualidade depende também das relações entre entradas e saída, da função
de perda e da cobertura das condições operacionais.

O escalonamento das variáveis de entrada pode reduzir diferenças de
magnitude e facilitar o ajuste numérico. Duas formas usuais são a
padronização \textit{z-score},

\begin{equation}
\label{eq:zscore}
z=\frac{x-\mu}{\sigma},
\end{equation}

\noindent em que $\mu$ e $\sigma$ são, respectivamente, média e desvio
padrão, e a normalização \textit{min-max},

\begin{equation}
\label{eq:minmax}
x'=\frac{x-x_{\min}}{x_{\max}-x_{\min}},
\end{equation}

\noindent em que $x_{\min}$ e $x_{\max}$ delimitam a escala de
referência. A primeira produz variáveis centradas e expressas em desvios
padrão; a segunda mapeia para $[0,1]$ as entradas contidas no intervalo
adotado, desde que $x_{\max}\neq x_{\min}$. Os parâmetros de ambas devem
ser estimados somente com o conjunto de treinamento, evitando vazamento
de informação. A padronização requer $\sigma>0$; variáveis constantes
exigem tratamento específico para evitar divisões por zero.

Para tratar a assimetria de um alvo contínuo, a família de transformações
de Box--Cox \cite{boxcox1964} é definida para argumentos positivos.
Quando existem observações nulas, admite-se um deslocamento
$\varepsilon>0$ tal que $u=y+\varepsilon>0$. A variável transformada é

\begin{equation}
\label{eq:boxcox}
z=g_{\lambda}(y+\varepsilon)=
\begin{cases}
\displaystyle\frac{(y+\varepsilon)^{\lambda}-1}{\lambda}, & \lambda\neq 0,\\[4pt]
\ln(y+\varepsilon), & \lambda=0.
\end{cases}
\end{equation}

A recuperação do alvo original exige a transformação inversa e a
subtração do deslocamento:

\begin{equation}
\label{eq:boxcox_inversa}
\widehat{y}=g_{\lambda}^{-1}(\widehat{z})-\varepsilon=
\begin{cases}
(1+\lambda\widehat{z})^{1/\lambda}-\varepsilon,
& \lambda\neq 0,\\
\exp(\widehat{z})-\varepsilon, & \lambda=0.
\end{cases}
\end{equation}

O parâmetro $\lambda$ controla a forma da transformação e costuma ser
estimado por máxima verossimilhança no conjunto de treinamento. Valores
menores que 1 tendem a comprimir a cauda superior e aproximar a
distribuição de uma normal; valores maiores que 1 atuam na direção
oposta. No caso $\lambda\neq0$, a inversão requer
$1+\lambda\widehat{z}>0$, o que impõe um limite ao domínio da predição. A
transformação altera a escala de treinamento e pode modificar a
assimetria e a dispersão, mas não garante, isoladamente, melhoria
preditiva na unidade física. Seu deslocamento e o parâmetro $\lambda$
devem ser definidos com base apenas nos dados de treinamento.

Neste trabalho, a transformação é aplicada ao alvo do MLP, não às
entradas, e a predição é convertida de volta para kW antes do cálculo das
métricas. A regressão linear e o RF são ajustados diretamente em kW; a
comparação entre os modelos, portanto, ocorre na unidade física, mas o
MLP aprende em um domínio transformado e pode apresentar diferenças de
comportamento nas faixas de potência mais altas.

\section{Modelos de regressão para estimação de potência}
\label{sec:ft-modelos}

\subsection{Regressão linear múltipla}
\label{subsec:ft-regressao-linear}

A regressão linear múltipla representa a saída como uma combinação afim
das entradas:

\begin{equation}
\label{eq:ft-regressao-linear}
\hat{y}=\beta_0+\sum_{j=1}^{p}\beta_j x_j,
\end{equation}

\noindent em que $\beta_0$ é o intercepto e $\beta_j$ são os coeficientes
ajustados. Seus coeficientes podem ser estimados por mínimos quadrados
ordinários, que minimizam a soma dos erros ao quadrado. A baixa
complexidade e a interpretabilidade dos coeficientes fazem dela uma
referência natural para investigar se modelos não lineares oferecem
ganhos que justifiquem maior custo computacional. A forma linear, no
entanto, não representa a saturação do inversor nem as interações entre
irradiância e temperatura, o que limita seu desempenho nos extremos da
faixa de operação.

\subsection{Random Forest}
\label{subsec:ft-rf}

O \textit{Random Forest} combina árvores de decisão, cuja média fornece
a predição em problemas de regressão \cite{breiman2001}:

\begin{equation}
\label{eq:ft-rf}
\hat{y}(\mathbf{x})=\frac{1}{B}\sum_{b=1}^{B}T_b(\mathbf{x}).
\end{equation}

Cada árvore é treinada em uma amostra \textit{bootstrap} do conjunto de
treinamento e, em cada divisão, considera apenas um subconjunto
aleatório das entradas. Essas duas fontes de aleatoriedade reduzem a
correlação entre as árvores e tornam a média mais estável do que uma
única árvore. As árvores individuais podem ser profundas e capturar
interações não lineares, enquanto a agregação controla a variância.

O número de árvores, o mínimo de amostras por folha e a profundidade
máxima controlam o compromisso entre ajuste e complexidade
\cite{probst2018}. Florestas
maiores tendem a estabilizar a predição, mas aumentam o tempo de
treinamento e o tamanho da implementação embarcada; folhas menores
permitem ajustes mais finos e também favorecem o sobreajuste. Embora
competitivos em dados tabulares \cite{grinsztajn2022}, os métodos
baseados em árvores não são universalmente superiores às redes neurais,
e a comparação deve ser feita nas mesmas condições de dados e de
avaliação.

Para implantação, a floresta pode ser convertida em código C, com as
divisões e os valores das folhas embutidos no firmware. O custo depende
do número de nós e de folhas, não de uma multiplicação de matrizes; por
isso, o tamanho da floresta precisa ser tratado como parte do projeto, e
não apenas como detalhe de treinamento.

\subsection{Perceptron de múltiplas camadas}
\label{subsec:ft-mlp}

O perceptron de múltiplas camadas (\textit{Multilayer Perceptron}, MLP)
combina transformações lineares e ativações não lineares. Para uma
camada $\ell$, tem-se

\begin{equation}
\label{eq:ft-mlp}
\mathbf{h}^{(\ell)}=\phi\left(
\mathbf{W}^{(\ell)}\mathbf{h}^{(\ell-1)}+\mathbf{b}^{(\ell)}
\right),
\end{equation}

\noindent em que $\mathbf{W}^{(\ell)}$ e $\mathbf{b}^{(\ell)}$ são os
parâmetros aprendidos e $\phi$ é a ativação. A função ReLU é dada por
$\phi(a)=\max(0,a)$ e utiliza operações simples, compatíveis com a
execução quantizada. Em comparação com ativações saturantes, a ReLU
mantém o gradiente não nulo para entradas positivas, o que facilita o
treinamento de redes com várias camadas \cite{glorot2011}. A camada de saída é linear em
problemas de regressão, de modo que a rede produz uma estimativa
contínua em vez de uma probabilidade.

A profundidade e a largura das camadas ocultas determinam a capacidade
do modelo. Com unidades suficientes, a rede pode aproximar funções
contínuas \cite{hornik1989}, mas a escolha da arquitetura envolve um compromisso entre
capacidade, tempo de treinamento e custo de inferência. Para limitar o
sobreajuste, podem ser empregados \textit{dropout}, que desativa
unidades durante o treinamento, e parada antecipada com base no erro de
validação \cite{prechelt1998,srivastava2014}. A parada antecipada
restaura os parâmetros do melhor ponto observado, e não os da última
época.

A perda de Huber combina comportamento quadrático para erros pequenos e
linear para erros elevados \cite{huber1964}. Definindo o resíduo
$e=y-\hat y$ e o limiar $\delta>0$, tem-se

\begin{equation}
\label{eq:ft-huber}
L_{\delta}(e)=
\begin{cases}
\dfrac{e^2}{2}, & |e|\leq\delta,\\[4pt]
\delta\left(|e|-\dfrac{\delta}{2}\right), & |e|>\delta.
\end{cases}
\end{equation}

O comportamento linear reduz a influência de observações atípicas,
enquanto a região quadrática mantém a suavidade próxima de zero. O
limiar controla o ponto de transição e precisa ser definido em uma
escala compatível com o alvo.

O otimizador Adam emprega estimativas adaptativas dos momentos do
gradiente para atualizar os parâmetros durante o treinamento
\cite{kingma2015}. A taxa de aprendizado efetiva é ajustada por
estimativas de primeira e segunda ordem, com correção de viés nas
primeiras iterações. A função de perda, o limiar e as configurações de
otimização são escolhas metodológicas que devem ser especificadas para
cada experimento e validadas sem consultar o conjunto de teste.

\section{Avaliação e calibração dos estimadores}
\label{sec:ft-avaliacao}

\subsection{Métricas e interpretação dos erros}
\label{subsec:ft-metricas}

Em regressão, sejam $y_i$ o valor medido, $\hat y_i$ o valor estimado,
$\bar y$ a média das medições e $n$ o número de pares válidos. O erro
absoluto médio (MAE) e a raiz do erro quadrático médio (RMSE) são
definidos respectivamente por

\begin{equation}
\label{eq:mae}
\mathrm{MAE}=\frac{1}{n}\sum_{i=1}^{n}|y_i-\hat y_i|,
\end{equation}

\begin{equation}
\label{eq:rmse}
\mathrm{RMSE}=\sqrt{\frac{1}{n}\sum_{i=1}^{n}(y_i-\hat y_i)^2}.
\end{equation}

O MAE mantém a unidade da potência e trata todos os erros com o mesmo
peso; o RMSE atribui maior peso aos desvios elevados, porque eleva o
erro ao quadrado antes de calcular a média. A comparação entre os dois
ajuda a identificar se o modelo concentra erros em poucas observações ou
se erra de forma uniforme. O coeficiente de determinação compara o erro
do modelo com a variabilidade do conjunto de referência:

\begin{equation}
\label{eq:r2}
R^2=1-\frac{\sum_{i=1}^{n}(y_i-\hat y_i)^2}
{\sum_{i=1}^{n}(y_i-\bar y)^2}.
\end{equation}

Seu denominador exige variabilidade não nula de $y_i$. Valores próximos
de 1 indicam que o modelo explica a maior parte da variância; valores
próximos de 0 indicam desempenho semelhante ao de uma média constante;
valores negativos são possíveis quando o modelo é pior que essa
referência. Além disso, valores de $R^2$ de conjuntos com dispersões
distintas não são diretamente equivalentes, o que impede comparar
resoluções diferentes apenas por esse indicador.

O erro percentual absoluto médio (MAPE) é expresso por

\begin{equation}
\label{eq:mape}
\mathrm{MAPE}=\frac{100}{n_{+}}
\sum_{i:y_i>0}\left|\frac{y_i-\hat y_i}{y_i}\right|,
\end{equation}

\noindent em que $n_{+}$ é a quantidade de observações com referência
estritamente positiva. Essa convenção exclui denominadores nulos, mas
não remove a instabilidade do MAPE para potências muito pequenas
\cite{hyndman2006}. O viés médio, com sinal positivo para superestimação,
é definido por

\begin{equation}
\label{eq:vies}
\text{Viés}=\frac{1}{n}
\sum_{i=1}^{n}(\hat y_i-y_i).
\end{equation}

O viés pode se aproximar de zero mesmo na presença de erros individuais
relevantes, porque compensações entre superestimativas e subestimativas
se anulam na média. O percentil 90 do erro absoluto (P90) corresponde ao
quantil de ordem $0{,}90$ dos valores $|y_i-\hat y_i|$. Diferentemente da
média, ele descreve a cauda da distribuição de erros e é menos sensível a
um único valor extremo do que a média. A interpretação conjunta dessas métricas e de
gráficos de valores medidos versus estimados é mais informativa que a
análise de um único indicador \cite{chicco2021,hodson2022}. Nos
resultados, a comparação entre granularidades combina MAE, RMSE, $R^2$,
viés e P90, porque cada um responde a uma pergunta diferente.

\subsection{Calibração de viés e mudanças de distribuição}
\label{subsec:ft-calibracao}

Uma correção aditiva de viés utiliza observações de calibração para
estimar o desvio médio entre a grandeza medida e a predita. Para um
conjunto com $n_{\mathrm{cal}}$ pares válidos, pode-se definir

\begin{equation}
\label{eq:offset_calibracao}
c_{\mathrm{cal}}=\frac{1}{n_{\mathrm{cal}}}
\sum_{i=1}^{n_{\mathrm{cal}}}(y_i-\hat y_i),
\qquad
\hat y_{i,\mathrm{corr}}=\hat y_i+c_{\mathrm{cal}}.
\end{equation}

O sinal positivo de $c_{\mathrm{cal}}$ aumenta a estimativa, enquanto o
negativo a reduz. Esse ajuste pode atenuar o erro sistemático, mas não
elimina a dispersão residual nem garante desempenho fora das condições
da calibração \cite{koo2022calibracao}. A correção é uma constante e,
por isso, não corrige viés condicional: ela desloca toda a distribuição
das estimativas e mantém a forma dos erros. Se o modelo subestima uma
faixa de potência e superestima outra, o valor médio pode ficar próximo
de zero mesmo com erros relevantes em ambos os extremos. A análise por
faixa de potência ou por período é necessária para detectar esse tipo de
comportamento.

Mudanças sazonais, operacionais ou entre inversores podem alterar a
relação aprendida, exigindo avaliação em períodos e equipamentos
distintos \cite{zhang2023drift}. Por isso, a janela de calibração deve ser
independente do teste e representativa das condições em que o estimador
será usado. No presente trabalho, a calibração utiliza o segundo
semestre de 2024 e o teste, o ano de 2025, sem participação das
observações de teste no cálculo do \textit{offset}.

\section{Inferência em sistemas embarcados}
\label{sec:ft-embarcado}

\subsection{TinyML e quantização pós-treinamento}
\label{subsec:ft-tinyml}

TinyML compreende a execução de modelos de aprendizado de máquina em
dispositivos de recursos limitados. Implementações como TensorFlow Lite
Micro permitem inferência em microcontroladores sem sistema operacional
convencional \cite{david2021}. O custo de implantação inclui não apenas
o arquivo de parâmetros, mas também código, buffers e memória de
trabalho. Por isso, a escolha do modelo não pode considerar apenas a
precisão obtida no computador; ela precisa incluir o consumo de memória
e o tempo de inferência no hardware alvo.

A quantização pós-treinamento (\textit{Post-Training Quantization}, PTQ)
converte um modelo previamente treinado em ponto flutuante para uma
representação de precisão reduzida, sem novo treinamento. Diferentemente,
a \textit{Quantization-Aware Training} (QAT) incorpora os efeitos da
quantização ao processo de treinamento ou ajuste fino
\cite{krishnamoorthi2018quantization,nagel2021quantization}. A PTQ é
mais simples de aplicar, porque parte do modelo já validado, mas pode
introduzir perdas maiores quando a quantização é agressiva ou quando a
distribuição das ativações é assimétrica.

Na modalidade de quantização inteira de 8 bits (\textit{full-integer
quantization}), pesos e ativações do modelo são representados por
inteiros. Uma relação afim aproxima o valor real por \cite{jacob2018}:

\begin{equation}
\label{eq:ft-quantizacao}
x\approx s(q-z),
\end{equation}

\noindent em que $q$ é o valor inteiro, $s$ é a escala e $z$ é o ponto
zero. A calibração da PTQ utiliza amostras representativas para estimar
os parâmetros de quantização das ativações; não se confunde com a
normalização estatística das entradas nem com a correção de viés das
predições. A redução do tamanho do modelo pode introduzir diferenças
numéricas e exige comparar precisão, memória e tempo de inferência. Neste trabalho, a PTQ é aplicada ao MLP já
treinado, e a variante inteira é comparada com a versão em ponto
flutuante no mesmo conjunto de teste. A comparação permite separar a
perda atribuída à precisão numérica da diferença entre arquiteturas,
porque as duas variantes mantêm a mesma estrutura de rede.

\subsection{Tempo de inferência e limites da avaliação embarcada}
\label{subsec:ft-latencia}

O tempo de inferência deve ser distinguido da latência de comunicação,
da inicialização e do processamento de lotes. Medidas por amostra, com
média e percentis, permitem caracterizar o comportamento do firmware,
inclusive eventuais custos de aquecimento. A média bruta pode ser
influenciada pelas primeiras execuções, que pagam custos de cache,
inicialização de buffers ou ajuste de frequência; por isso, a média sem
aquecimento e o P90 complementam a média geral.

A execução com reprodução de dados históricos comprova o funcionamento
do modelo no hardware e permite quantificar sua precisão nesse cenário,
mas não substitui a validação em campo com aquisição contínua e falhas
reais de instrumentação. O tempo medido também depende do modelo, do
runtime e da frequência do processador, o que impede comparar plataformas
distintas sem controlar essas condições. A avaliação embarcada deve,
portanto, ser lida como uma etapa intermediária entre o teste no
computador e a operação em campo.
